% ABSTRACT — draft 2026-10-26. JASA limit 200 words. Bracketed bold text is a
% placeholder until data collection closes (13.11.2026).
Listeners can learn new spectral cues for sound elevation, but it is not known
whether the learned spectral-to-spatial mapping belongs to the ear that was
trained or is available to both ears. We used a virtual auditory environment to
replace the fine spectral detail of one ear's head-related transfer functions
with that of another listener, leaving the spectral envelope and the interaural
time and level differences intact, while the other ear received only its own
spectral envelope. Participants trained with this modified cue on four
consecutive days. Because the cue was rendered rather than worn, the identical
filter could then be delivered to the untrained ear. On the final day,
localization was tested in a $2 \times 2$ design crossing the ear that carried
the modified cue (trained, untrained) with the side of space (trained,
mirrored). \textbf{[RESULT: adaptation with the trained ear --- elevation gain
X $\to$ Y, polar error $-Z^\circ$.] [RESULT: change with the identical cue on
the untrained ear, same metrics.] [INTERPRETATION: ear-specific vs shared
mapping, one sentence.]}
